\chapter*{Parte V. BPE paso a paso y buenas prácticas de comparación}
\addcontentsline{toc}{chapter}{Parte V. BPE y validación de modelos}
\section*{29. Cinco palabras, una regla de fusión}
La técnica de BPE puede ilustrarse sin emplear bibliotecas de aprendizaje profundo. Usaremos las palabras \emph{casa} (5 apariciones), \emph{casas} (2), \emph{cosa} (6), \emph{cosas} (3) y \emph{caso} (4). Cada palabra se representa inicialmente como caracteres y un marcador de fin $\langle/\mathrm{w}\rangle$. La frecuencia de una palabra pondera todas sus parejas adyacentes. El objetivo de este ejercicio no es crear un tokenizador listo para producción, sino entender qué se modifica al aprender unidades.

\begin{center}
\begin{tabular}{lcl}\toprule
Palabra & Frecuencia & Segmentación inicial\\\midrule
casa &5& c $\mid$ a $\mid$ s $\mid$ a $\mid$ $\langle/\mathrm{w}\rangle$\\
casas &2& c $\mid$ a $\mid$ s $\mid$ a $\mid$ s $\mid$ $\langle/\mathrm{w}\rangle$\\
cosa &6& c $\mid$ o $\mid$ s $\mid$ a $\mid$ $\langle/\mathrm{w}\rangle$\\
cosas &3& c $\mid$ o $\mid$ s $\mid$ a $\mid$ s $\mid$ $\langle/\mathrm{w}\rangle$\\
caso &4& c $\mid$ a $\mid$ s $\mid$ o $\mid$ $\langle/\mathrm{w}\rangle$\\\bottomrule
\end{tabular}
\end{center}
\noindent En este corpus de juguete, el par de caracteres $(a,s)$ aparece ponderadamente $5+2\cdot2+3+4=16$ veces (en \emph{casas} aparece dos veces). El par $(s,a)$ también suma 16: hay un empate que resolvemos con la regla lexicográfica declarada. Se elige $(a,s)$ y se fusiona en una unidad \texttt{as}. En \emph{casa} pasamos de \texttt{c|a|s|a|fin} a \texttt{c|as|a|fin}. El cambio no altera el texto original; altera la forma de \emph{segmentarlo}.

\section*{30. Cómo leer una tabla de fusiones}
Las seis primeras fusiones, bajo conteo ponderado, empates lexicográficos y sustitución de pares no solapados, resultan así:
\begin{center}
\begin{tabular}{c l r r}\toprule
Paso & Par fusionado & Frecuencia & Unidades después\\\midrule
0 & -- & -- &105\\
1 & $(a,s)$ &16&89\\
2 & $(a,\langle/\mathrm{w}\rangle)$ &11&78\\
3 & $(c,as)$ &11&67\\
4 & $(c,o)$ &9&58\\
5 & $(co,s)$ &9&49\\
6 & $(cos,a\langle/\mathrm{w}\rangle)$ &6&43\\\bottomrule
\end{tabular}
\end{center}
\noindent Cada resta $105-89=16$, $89-78=11$, etc. es verificable en el módulo R/Python. No debemos suponer sin comprobar que \emph{frecuencia de un par} y \emph{sustituciones efectivas} siempre coinciden: en \texttt{aaaaa}, el par $(a,a)$ aparece en cuatro ventanas superpuestas, pero una pasada de fusión no solapada puede efectuar solamente dos reemplazos.

\section*{31. Una tabla de preguntas para discutir con la clase}
\begin{center}
\begin{tabularx}{\textwidth}{p{55mm}X}\toprule
Pregunta & Idea matemática o metodológica\\\midrule
¿BPE siempre disminuye el número de unidades? & Una fusión aplicada sustituye dos unidades por una; las reducciones efectivas dependen de los solapamientos.\\
¿BPE produce mejores predicciones? & No se deduce de reducir longitud: es necesario evaluar modelos probabilísticos en prueba.\\
¿Podemos comparar PPL entre segmentaciones? & No directamente por token; fijar texto, unidades y probabilidades de la cadena completa.\\
¿El corpus técnico permite generalizar al español técnico? & Solo con muestreo adecuado, procedencia y dominio especificados.\\\bottomrule
\end{tabularx}
\end{center}
\textbf{Puente hacia el proyecto IC10.} Podemos investigar si ciertos términos técnicos ocupan más tokens por palabra que la prosa general. Esto ya es una medición falsable y puede ampliarse a BPC con modelos adecuadamente controlados. La tabla del catálogo identifica D01 y D02 como corpus generales candidatos; ninguna cifra empírica sobre esos corpus se afirma en esta guía.
\begin{idea}[Tres capas: prompt, contexto y arnés]
En un modelo generativo, el \emph{prompt} aporta la entrada inicial; el \emph{contexto} es la historia efectiva disponible, según la ventana y las reglas de truncamiento; el \emph{arnés} ejecuta y registra el procedimiento (muestreo, parada, validación, semillas y herramientas). Que los tres influyan en la salida no los convierte en la misma operación. Su estudio sistemático se retomará en sesiones posteriores.
\end{idea}
\clearpage
